\documentclass[11pt]{article}

% -------------------- Packages --------------------
\usepackage[a4paper,margin=1in]{geometry}
\usepackage{hyperref}
\usepackage{enumitem}
\usepackage{titlesec}
\usepackage{fancyhdr}
\usepackage{xcolor}
\usepackage{listings}
% listings has no built-in yaml or Docker language; define them (plain, no highlighting) so latexmk builds without errors.
\lstdefinelanguage{yaml}{}
\lstdefinelanguage{Docker}{}
\usepackage{courier}

% -------------------- Hyperref Setup --------------------
\hypersetup{
  colorlinks=true,
  linkcolor=blue,
  urlcolor=blue,
  citecolor=blue
}

% -------------------- Header/Footer --------------------
\pagestyle{fancy}
\fancyhf{}
\lhead{MCV Lab}
\rhead{Containers \& Orchestration}
\cfoot{\thepage}

% -------------------- Section Formatting --------------------
\titleformat{\section}{\Large\bfseries}{\thesection}{0.8em}{}
\titleformat{\subsection}{\normalsize\bfseries}{\thesubsection}{0.8em}{}
\titleformat{\subsubsection}{\normalsize\bfseries}{\thesubsubsection}{0.8em}{}

\titlespacing*{\section}{0pt}{1.5em}{0.8em}
\titlespacing*{\subsection}{0pt}{1.0em}{0.5em}
\titlespacing*{\subsubsection}{0pt}{0.8em}{0.4em}

% -------------------- Global List Formatting --------------------
\setlist[itemize]{
  label={-},
  leftmargin=2em,
  labelsep=0.6em,
  itemsep=0.2em
}
\setlist[itemize,2]{itemsep=0.15em}
\setlist[enumerate]{leftmargin=2em,itemsep=0.15em}

% -------------------- Listings / Code Style --------------------
\definecolor{codegray}{rgb}{0.95,0.95,0.95}
\definecolor{commentgreen}{rgb}{0.0,0.5,0.0}

\lstset{
  basicstyle=\footnotesize\ttfamily,
  backgroundcolor=\color{codegray},
  frame=single,
  breaklines=true,
  showstringspaces=false,
  keywordstyle=\color{blue},
  commentstyle=\color{commentgreen},
  stringstyle=\color{red},
  columns=fullflexible,
  keepspaces=true,
  tabsize=2,
  xleftmargin=0.02\textwidth,
  xrightmargin=0.02\textwidth
}

% -------------------- Document Metadata --------------------
\title{\textbf{Medical Computer Vision Lab \\ \small{
Docker vs Kubernetes: Lab Standard for Reproducible ML \& Scalable Execution}}}
\author{Maintainer: Ashery Mbilinyi}
\date{Version 1.0 : \today}

\begin{document}
\maketitle

\tableofcontents
\newpage

\section{Purpose}
This document explains the difference between \textbf{Docker} and \textbf{Kubernetes}, and defines a practical lab standard for using containers in research workflows.
The goal is to ensure our ML work is:
\begin{itemize}
  \item \textbf{Reproducible}: the same environment runs across laptops, servers, and cloud
  \item \textbf{Portable}: experiments can move between machines with minimal friction
  \item \textbf{Scalable}: training can be executed reliably on HPC schedulers (e.g., SLURM) and, when needed, on clusters
  \item \textbf{Collaborative}: new lab members can run baselines quickly
\end{itemize}

\section{Key Definitions}

\subsection{What is a Container?}
A \textbf{container} packages:
\begin{itemize}
  \item code + dependencies (Python libraries, OS packages)
  \item runtime settings (environment variables, entrypoints)
  \item a filesystem snapshot (image layers)
\end{itemize}
This makes training and inference behave consistently across machines.

\subsection{Docker (Containerization)}
\textbf{Docker} is primarily a toolchain to \textbf{build} and \textbf{run} containers.

\textbf{What Docker gives us:}
\begin{itemize}
  \item \textbf{Dockerfile} $\rightarrow$ builds a \textbf{Docker image}
  \item \textbf{docker run} $\rightarrow$ runs the image as a container
  \item reproducible dependency management (CUDA/PyTorch stacks, system libraries)
\end{itemize}

\subsection{Kubernetes (Orchestration)}
\textbf{Kubernetes (K8s)} is a \textbf{container orchestration system}. It schedules and manages containers across a \textbf{cluster} of machines.

\textbf{What Kubernetes adds on top of containers:}
\begin{itemize}
  \item automatic scheduling of workloads across many nodes
  \item scaling (up/down) and self-healing (restart/reschedule)
  \item service discovery/networking between components
  \item rolling updates and deployment management
\end{itemize}

\section{Docker vs Kubernetes: Practical Differences}

\subsection{Mental Model}
\begin{itemize}
  \item \textbf{Docker} = package and run the application reliably
  \item \textbf{Kubernetes} = manage many running containers reliably at cluster scale
\end{itemize}

\subsection{Comparison Table}
\begin{itemize}
  \item \textbf{Docker} is typically enough for: local dev, single GPU server runs, reproducible environments.
  \item \textbf{Kubernetes} is useful for: cluster-wide scheduling, production deployments, multi-service platforms.
\end{itemize}

\section{Lab Standard (Recommended): Docker + SLURM (Academic Default)}
In most academic environments, the standard scaling layer is \textbf{SLURM}, not Kubernetes.

\subsection{Why Docker + SLURM is the Default}
\begin{itemize}
  \item Docker gives reproducible environments.
  \item SLURM schedules jobs on HPC resources (GPUs/CPUs) and handles fair-sharing.
  \item This is simpler to operate than Kubernetes for typical lab needs.
\end{itemize}

\subsection{Important Note: Docker vs Apptainer/Singularity on HPC}
Many HPC systems disallow Docker for security reasons.
Instead, they use \textbf{Apptainer/Singularity} to run containerized workloads.

\textbf{Practical pattern:}
\begin{lstlisting}[language=bash]
Docker image  -->  Apptainer/Singularity image  -->  Run via SLURM
\end{lstlisting}

\section{When We Use Docker}

\subsection{Required (Lab Policy)}
All ML repositories should provide a container recipe:
\begin{itemize}
  \item \texttt{Dockerfile} (preferred), or
  \item Apptainer recipe / documented conversion steps (if HPC requires it)
\end{itemize}

\subsection{Typical Docker Commands}
\begin{lstlisting}[language=bash]
# Build
docker build -t mcv-project:latest .

# Run (GPU)
docker run --gpus all -it --rm \
  -v $(pwd):/workspace \
  mcv-project:latest
\end{lstlisting}

\subsection{What Belongs in the Docker Image}
\begin{itemize}
  \item pinned Python dependencies (requirements.txt / poetry.lock / conda env)
  \item system libs needed for the project (e.g., OpenCV deps)
  \item fixed CUDA/PyTorch compatibility (base image choice matters)
\end{itemize}

\section{When Kubernetes is Worth It (Optional)}
Kubernetes is not required for most lab research workflows. It becomes valuable if we run a shared platform, for example:
\begin{itemize}
  \item multiple groups share a GPU cluster and need a common scheduling layer
  \item we host persistent services (APIs, dashboards, model serving endpoints)
  \item we run many pipelines continuously (training + evaluation + deployment)
\end{itemize}

\subsection{Kubernetes Concepts (Minimal Glossary)}
\begin{itemize}
  \item \textbf{Pod}: smallest deployable unit (one or more containers)
  \item \textbf{Deployment}: manages a replicated set of pods (rolling updates)
  \item \textbf{Job}: runs a pod to completion (batch workloads)
  \item \textbf{Service}: stable network endpoint for accessing pods
  \item \textbf{Ingress}: routes external traffic to services
\end{itemize}

\section{Recommended Repository Structure (Containers + Experiments)}
All projects should follow a consistent structure:
\begin{lstlisting}
project/
  Dockerfile
  requirements.txt (or environment.yml)
  src/
  scripts/
    train.py
    eval.py
  configs/
    baseline.yaml
  slurm/
    train.slurm
  runs/               # outputs (gitignored)
  README.md
\end{lstlisting}

\section{Best Practices}

\subsection{Reproducibility}
\begin{itemize}
  \item Pin dependencies (exact versions).
  \item Save runtime config per run (\texttt{config.json/yaml} + \texttt{cmd.txt}).
  \item Save checkpoints (best + last) and log metrics.
\end{itemize}

\subsection{Security and Practical Notes}
\begin{itemize}
  \item Do not bake secrets (API keys) into images.
  \item Prefer environment variables or secret managers.
  \item Use read-only mounts where possible when evaluating.
\end{itemize}

\section{Checklist (Code Review)}
Before merging a new training pipeline:
\begin{itemize}
  \item \textbf{Container exists}: repo includes \texttt{Dockerfile} (or Apptainer recipe).
  \item \textbf{One-command run}: README has build + run instructions.
  \item \textbf{Reproducible outputs}: runs save config + command + checkpoints.
  \item \textbf{HPC compatibility}: provide \texttt{slurm/train.slurm} if used on cluster.
\end{itemize}

\section{Appendix: Minimal Templates}

\subsection{Minimal Dockerfile (PyTorch + CUDA)}
\begin{lstlisting}[language=Docker]
FROM pytorch/pytorch:2.3.1-cuda12.1-cudnn8-runtime

WORKDIR /workspace

# System deps (adjust as needed)
RUN apt-get update && apt-get install -y --no-install-recommends \
    git curl ca-certificates \
    && rm -rf /var/lib/apt/lists/*

COPY requirements.txt /workspace/requirements.txt
RUN pip install --no-cache-dir -r requirements.txt

COPY . /workspace

# Default command (override in docker run if needed)
CMD ["python", "scripts/train.py", "--help"]
\end{lstlisting}

\subsection{Minimal SLURM Script (Runs Containerized Training)}
\begin{lstlisting}[language=bash]
#!/bin/bash
#SBATCH --job-name=mcv_train
#SBATCH --gres=gpu:1
#SBATCH --cpus-per-task=8
#SBATCH --mem=32G
#SBATCH --time=12:00:00
#SBATCH --output=slurm-%j.out

module load apptainer

# Example: run an Apptainer image built from Docker
apptainer exec --nv mcv-project.sif \
  python scripts/train.py --config configs/baseline.yaml
\end{lstlisting}

\end{document}
